El análisis parte de la tabla \texttt{PSCompPars} del NASA Exoplanet Archive, específicamente del archivo compacto \texttt{data/PSCompPars\_2026.04.25\_17.36.36.csv}. Esta versión contiene 6273 exoplanetas y 84 columnas. A partir de esta tabla se construyó una versión procesada e imputada del catálogo para poder alimentar los espacios de variables usados por Mapper.

La imputación fue necesaria porque Mapper requiere una matriz numérica completa para construir los grafos. Si se trabajara únicamente con casos completos, se perdería una parte importante del catálogo y algunos espacios de análisis quedarían severamente reducidos. El problema es especialmente visible en las variables térmicas, como \texttt{pl\_insol} y \texttt{pl\_eqt}, y en variables que no estaban completas en la tabla compacta, como \texttt{pl\_dens}. Por esta razón, el pipeline no solo imputó valores faltantes, sino que también conservó trazabilidad sobre qué valores fueron observados, derivados físicamente o imputados estadísticamente.

\subsection*{Disponibilidad inicial de variables}

Las variables principales usadas para Mapper presentaban distintos niveles de faltantes antes de la imputación o derivación. La Tabla~\ref{tab:raw_missing_values} resume la disponibilidad inicial de las variables principales.

\begin{table}[H]
\centering
\small
\begin{tabular}{lrrp{0.38\textwidth}}
\toprule
\textbf{Variable} & \textbf{Faltantes crudos} & \textbf{Total} & \textbf{Comentario} \\
\midrule
\texttt{pl\_rade} & 50 & 6273 & Radio planetario; variable física principal. \\
\texttt{pl\_bmasse} & 31 & 6273 & Masa planetaria; variable física principal. \\
\texttt{pl\_dens} & 6273 & 6273 & No venía completa en la tabla compacta; se derivó para la mayoría de planetas. \\
\texttt{pl\_orbper} & 339 & 6273 & Periodo orbital; variable orbital principal. \\
\texttt{pl\_orbsmax} & 428 & 6273 & Semieje mayor; se derivó cuando fue posible. \\
\texttt{pl\_insol} & 1876 & 6273 & Insolación; variable térmica con alta incompletitud. \\
\texttt{pl\_eqt} & 1600 & 6273 & Temperatura de equilibrio; variable térmica con alta incompletitud. \\
\bottomrule
\end{tabular}
\caption{Valores faltantes crudos en las variables principales antes de imputación o derivación física.}
\label{tab:raw_missing_values}
\end{table}

Esta tabla muestra que no todas las variables tienen el mismo nivel de confiabilidad inicial. Las variables físicas básicas, como radio y masa, tienen pocos faltantes. En cambio, las variables térmicas tienen una proporción mucho mayor de valores ausentes. Esto anticipa una diferencia importante en la interpretación posterior: un grafo puede ser topológicamente complejo, pero si esa complejidad depende de variables altamente imputadas, debe interpretarse con cautela.

\subsection*{Derivaciones físicas}

Antes de aplicar imputación estadística, el pipeline intentó completar algunas variables mediante relaciones físicas conocidas. Esta diferencia es importante: una variable derivada no es una observación directa, pero tampoco es una imputación puramente estadística. Es un valor calculado a partir de otras mediciones disponibles.

La primera derivación fue la densidad planetaria, calculada a partir de masa y radio:

\[
\texttt{pl\_dens}
=
5.514
\times
\frac{\texttt{pl\_bmasse}}{\texttt{pl\_rade}^{3}}.
\]

El factor 5.514 corresponde a la densidad media de la Tierra en g/cm$^3$, de modo que la fórmula convierte la relación masa-radio en una densidad planetaria aproximada. En la tabla compacta, \texttt{pl\_dens} tenía 6273 valores faltantes antes de la derivación. El pipeline derivó 6199 valores y dejó 74 faltantes después de esta etapa.

La segunda derivación fue el semieje mayor, usando una relación tipo Kepler a partir del periodo orbital y la masa estelar:

\[
\texttt{pl\_orbsmax}
=
\left(
\texttt{st\_mass}
\times
\left(
\frac{\texttt{pl\_orbper}}{365.25}
\right)^2
\right)^{1/3}.
\]

Antes de esta derivación, \texttt{pl\_orbsmax} tenía 428 valores faltantes. El pipeline derivó 420 valores y dejó 8 faltantes después de la etapa de derivación.

\begin{table}[H]
\centering
\small
\begin{tabular}{lrrrp{0.28\textwidth}}
\toprule
\textbf{Variable derivada} & \textbf{Faltantes antes} & \textbf{Derivados} & \textbf{Faltantes después} & \textbf{Fuente de derivación} \\
\midrule
\texttt{pl\_dens} & 6273 & 6199 & 74 & Masa y radio planetario. \\
\texttt{pl\_orbsmax} & 428 & 420 & 8 & Periodo orbital y masa estelar. \\
\bottomrule
\end{tabular}
\caption{Variables completadas mediante derivaciones físicas antes de la imputación estadística.}
\label{tab:physical_derivations}
\end{table}

Estas derivaciones tienen consecuencias directas para la interpretación. En particular, \texttt{pl\_dens} debe tratarse como una variable de control y no como evidencia completamente independiente cuando se usa junto con \texttt{pl\_rade} y \texttt{pl\_bmasse}. Como la densidad se calcula a partir de masa y radio, puede reforzar patrones ya contenidos algebraicamente en esas variables. Por eso el reporte compara explícitamente espacios con densidad y sin densidad, como \texttt{phys\_density} frente a \texttt{phys\_min}, y \texttt{joint} frente a \texttt{joint\_no\_density}.

\subsection*{Imputación estadística}

Después de las derivaciones físicas, el pipeline aplicó imputación estadística para completar la matriz de variables usada por Mapper. Se compararon tres métodos: \texttt{median}, \texttt{knn} e \texttt{iterative}. El método principal seleccionado fue \texttt{iterative}, porque obtuvo el mejor ranking promedio en las métricas de validación, considerando MAE y RMSE. En esta lectura, \texttt{median} funciona como una línea base robusta, \texttt{knn} como alternativa intermedia, e \texttt{iterative} como el método principal para los resultados visualizados e interpretados.

Es importante separar conceptualmente las variables derivadas de las imputadas. Una variable derivada se calcula usando una relación física y otras mediciones disponibles. Una variable imputada se estima estadísticamente a partir de patrones del dataset. Ambas introducen incertidumbre, pero no la misma clase de incertidumbre. Por eso el análisis conserva banderas de trazabilidad y no trata la matriz final como si todos sus valores tuvieran el mismo origen.

\begin{quote}
\textbf{Distinción metodológica.} En este reporte, ``derivado'' significa calculado mediante una relación física a partir de otras variables observadas; ``imputado'' significa estimado estadísticamente para completar valores faltantes. Ambos casos se auditan, pero no se interpretan como equivalentes.
\end{quote}

\subsection*{Trazabilidad de imputación y fuentes}

La imputación no se trató como un paso invisible del pipeline. Por planeta, la tabla completa imputada conservó columnas de trazabilidad que permiten identificar si un valor fue observado, faltante, derivado físicamente o imputado. Entre estas columnas se encuentran:

\[
\texttt{*\_was\_missing},\quad
\texttt{*\_was\_observed},\quad
\texttt{*\_was\_physically\_derived},\quad
\texttt{*\_was\_imputed},\quad
\texttt{*\_source},\quad
\texttt{imputation\_method}.
\]

Estas columnas aparecen especialmente en \texttt{reports/imputation/PSCompPars\_imputed\_iterative.csv}. Después, al construir los grafos Mapper, esta trazabilidad se agregó por nodo. Esto permitió calcular métricas de confianza como:

\[
\texttt{mean\_imputation\_fraction},\quad
\texttt{median\_imputation\_fraction},\quad
\texttt{observed\_fraction},\quad
\texttt{physically\_derived\_fraction},\quad
\texttt{imputed\_fraction},\quad
\texttt{frac\_any\_imputed}.
\]

Estas métricas aparecen en tablas como \texttt{node\_physical\_interpretation.csv} y \texttt{mapper\_node\_source\_audit.csv}. Su función es permitir que cada nodo o región del Mapper se evalúe no solo por su posición topológica, sino también por la calidad de la información que lo sostiene.

\subsection*{Confianza por espacio de variables}

La auditoría de imputación muestra que los espacios de variables no tienen el mismo nivel de confiabilidad. La Tabla~\ref{tab:imputation_by_space} resume las métricas principales para los grafos seleccionados.

\begin{table}[H]
\centering
\small
\begin{tabular}{lrrp{0.36\textwidth}}
\toprule
\textbf{Grafo seleccionado} & \textbf{Imputación media nodal} & \textbf{Riesgo} & \textbf{Interpretación} \\
\midrule
\texttt{phys\_min\_pca2\_cubes10\_overlap0p35} & 0.017 & Bajo & Espacio físico básico con poca dependencia de imputación. \\
\texttt{phys\_density\_pca2\_cubes10\_overlap0p35} & 0.002 & Bajo/moderado & Baja imputación, pero densidad frecuentemente derivada. \\
\texttt{orbital\_pca2\_cubes10\_overlap0p35} & 0.011 & Bajo & Espacio orbital confiable por imputación y candidato principal de inspección. \\
\texttt{joint\_no\_density\_pca2\_cubes10\_overlap0p35} & 0.176 & Moderado & Incluye variables térmicas, por lo que aumenta la dependencia de imputación. \\
\texttt{joint\_pca2\_cubes10\_overlap0p35} & 0.152 & Moderado & Espacio conjunto con densidad; combina riesgo térmico y redundancia de densidad. \\
\texttt{thermal\_pca2\_cubes10\_overlap0p35} & 0.588 & Alto & Estructura térmica con fuerte dependencia de imputación. \\
\bottomrule
\end{tabular}
\caption{Dependencia de imputación en los grafos seleccionados.}
\label{tab:imputation_by_space}
\end{table}

La lectura principal es que los espacios físico y orbital tienen baja dependencia de imputación, mientras que el espacio térmico tiene un riesgo mucho mayor. En particular, \texttt{thermal\_pca2\_cubes10\_overlap0p35} tiene una imputación media nodal de 0.588 y una fracción de nodos con alta imputación de 0.719. Esto no significa que el grafo térmico no contenga información, pero sí que cualquier interpretación astrofísica derivada de ese espacio debe ser más conservadora.

Los espacios conjuntos, \texttt{joint\_no\_density} y \texttt{joint}, tienen riesgo moderado porque incorporan variables térmicas como \texttt{pl\_insol} y \texttt{pl\_eqt}, que son precisamente las variables más incompletas. Por lo tanto, los espacios conjuntos son útiles para comparar perspectivas, pero no necesariamente producen evidencia más fuerte que los espacios simples. En este proyecto, la interpretación más confiable surge cuando una región combina baja imputación, coherencia física u orbital y baja asociación con sesgos de descubrimiento.

\subsection*{Implicaciones para la interpretación de Mapper}

La consecuencia metodológica de esta sección es que la imputación se usa como criterio de confianza. Un nodo Mapper no se interpreta únicamente por su posición en el grafo o por las variables promedio de sus planetas, sino también por el origen de sus datos. Una región con baja imputación puede sostener mejor una interpretación física u orbital. En cambio, una región con alta imputación, aunque sea topológicamente compleja, debe considerarse débil o cautelosa.

Esta distinción explica por qué el grafo orbital se vuelve especialmente importante en el reporte: combina baja imputación media nodal con alta complejidad topológica. En contraste, el grafo térmico puede mostrar estructura, pero su dependencia de imputación reduce la fuerza de cualquier conclusión física directa. Así, la imputación no aparece solo como un paso técnico previo, sino como una parte central de la evaluación científica de las regiones Mapper.